\documentclass[conference]{IEEEtran}
\usepackage{amsmath}
\usepackage{amssymb}
\usepackage{booktabs}
\usepackage{graphicx}
\usepackage{hyperref}
%%PREAMBLE%%
\usepackage{pdflscape}
\usepackage[margin=1in]{geometry}
% Wide result tables: landscape PDF pages via pdflscape.
\newenvironment{landscapetable}{%
 \clearpage
 \begin{landscape}%
}{%
 \end{landscape}%
 \clearpage
}

\title{Tables: Ant Colony Optimization for $k$-defective Maximum Edge Bicliques}
\author{Kiyaan Pillai}
\date{September 2026}

\begin{document}
\maketitle

Per-graph result tables for the paper. Section numbers below refer to that paper. Original graph sizes are in Appendix~\ref{app:sizes}. Each result table lists the reduced graph only.

\section{ACO vs $\theta$-heuristic}
The tables below list per-graph ACO results (100 ants, five epochs) against the $\theta$-heuristic. Table~\ref{tab:aco-heur-pn} is the primary paper setting ($k=2$, $\theta=5$); subsequent tables cover the $(k, \theta)$ sweep from Section~4.5. $U_k$ and $V_k$ are the reduced side sizes in thousands, and $|E_R|$ is the reduced edge count, after common-neighbor reduction.
Rows are grouped into three blocks separated by horizontal rules---ACO wins, $\theta$-heuristic wins, and ties or mutual $\theta$-infeasibility---using the win rule in Section~4. Within each block, rows are sorted by ACO $|E(D)|$.
For ACO, Time is the cumulative wall time through the winning replicate, and ETB is the epoch at which that replicate first reached its best $|E(D)|$. $|U_D|$, $|V_D|$, and $|E(D)|$ denote the returned biclique; all times are in milliseconds.
%%TABLE:k1t5i_PN%%
%%TABLE:k2t5i_PN%%
%%TABLE:k3t5i_PN%%
%%TABLE:k4t5i_PN%%
%%TABLE:k3t6i_PN%%
%%TABLE:k3t7i_PN%%

\section{Tug-of-War vs ACO and $\theta$-heuristic}
Table~\ref{tab:methods-k2t5i_heuristic-aco-tow} lists per-graph results for the Tug-of-War (TOW) alternating best-response heuristic against ACO and the $\theta$-heuristic. All three methods run on the same reduced graph. $U_k$ and $V_k$ are that graph's side sizes in thousands, and $|E_R|$ is its edge count. $|U_D|$, $|V_D|$, and $|E(D)|$ denote the returned biclique, and Time is the wall time of one timed run after a JIT warmup run, in milliseconds.
%%METHODS:k2t5i_heuristic-aco-tow%%

\section{Pivot seed comparison}
Table~\ref{tab:seed-compare-full} lists branch-and-pivot timings for every graph. $n_k$ is the reduced vertex count in thousands, and $|E_R|$ is the reduced edge count.
Rows are grouped into two blocks separated by a horizontal rule: graphs where ACO seeding reduced net end-to-end time (top), and graphs with no net savings or overhead (bottom). Within each block, rows are sorted by ascending absolute time saved.
Time saved is defined as $T_{\mathrm{pivot}}(\theta) - T_{\mathrm{discovery}} - T_{\mathrm{pivot}}(\mathrm{ACO})$, where $T_{\mathrm{discovery}}$ is the cost of running ACO for all 5 replicates. Time saved represents the net wall-clock time change of running branch-and-pivot with an ACO seed versus running branch-and-pivot on the $\theta$-heuristic seed alone. Positive values mean the ACO-seeded pipeline is faster overall. Negative values mean ACO search did not beat the $\theta$-heuristic, in which case $\text{time saved} = -T_{ACO}$. In some small graphs, ACO did beat the $\theta$-heuristic, but the time it took to run ACO meant it did not result in overall time savings. Reduction (\%) is defined as $100 \times \text{Time saved} / T_{\mathrm{pivot}}(\theta)$. Graphs where both ACO-seeded pivot and the $\theta$-seeded pivot are omitted, as defined in Section~3.3.
%%SEED_COMPARE%%

\appendix
\onecolumn
\section{Graph sizes}
\label{app:sizes}
%%TABLE:k2t5i_PN:sizes%%

\end{document}
